\section{Introdução}

Diante dos recentes avanços de modelos de IA (inteligência artifical), a Engenharia de Software vem se modificando para se adequar tais práticas. As ferramentas baseadas em IA tem auxiliado a apoiar a sociedade em tarefas cotidianas e do trabalho, dentro do ciclo de desenvolvimento, podemos apontar uma ajuda na licitação de requisitos, geração de código, testes e documentação.

De forma mais recente, agentes autônomos de programação tem se demonstrado capaz de implementar funcionalidades completas a partir de uma descrição usando linguagem natural, reduzindo assim o esforço manual necessário dos desenvolvedores de software.

Apesar dos avanços tecnológicos e da chegada de novos modelos poderosos, ainda é possível observar que a qualidade dos artefatos produzidos por agentes depende diretamente da qualidade das instruções fornecidas para eles. Requisitos incompletos, ambíguos e com uma estrutura errada podem gerar resultados inconsistentes durante o processo de desenvolvimento, aumentando o retrabalho e o número de interações necessárias para concluir tal tarefa. Por conta disso, é necessário utilizar técnicas avançadas de Engenharia de Prompt e refinamento de requisitos, para melhorar a comunicação entre desenvolvedores humanos e agentes baseados em Large Language Models (LLM).

Paralelo a isso, observa-se o surgimento de diversas abordagens diferentes para conseguirmos utilizar tais LLMs, dentre elas a utilização de múltiplos agentes, na qual cada um assume responsábilidade específica dentro do processo de desenvolvimento. Nesse cenário, um agente pode atuar com o refinamento dos requisitos enquanto outro é responsável pela implementação do código, reproduzindo de forma simplificada os papéis tradicionais que vemos nos profissionais do mercado de trabalho.

Este trabalho irá investigar uma abordagem de desenvolvimento orientada por agentes baseados na separação das atividades de refinamento de requisito e implementação de software. Como base para o estudo, foi desenvolvido um software para auxiliar a produção de textos acadêmicos utilizando \LaTeX{}, cujas funcionalidades e elementos de interface foram concebidos por meio de um processo de desenvolvimento assistido por agentes. A proposta consiste em utilizar um agente especializado em transformar as descrições e requisitos de funcionalidades escritas por humanos em prompts robustos que posteriormente serão utilizados como entrada para um agente autônomo de programação responsável por tal implementação de solução.

Como objetivo do trabalho, busca-se analisar o potencial dessa estratégia para apoiar o desenvolvimento de software assistido por IA. Serão discutidos seus benefícios, limitações e implicações para a Engenharia de Software. Para isso, foi conduzido um estudo de caso no qual os requisitos da aplicação forma inicialmente definidos por um desenvolvedor humano, refinados por um agente especializado e posteriormente implementados por um agente autônomo de programação.

A organização desse artigo está dividído em 4 seções, no qual a Seção 2 apresentaremos os conceitos fundamentais relacionados a agentes de software. A Seção 3 descreveremos a abordagem proposta e o estudo de caso realizado. Por fim, a Seção 4 apresenta as conclusões e perspectivas para trabalhos futuros.